% From Spectra to Topology: the Hodge--Dijkgraaf--Witten bridge.
% Theory primer for the tessera cobordism realizability work.
% Standalone whitepaper -- intentionally NOT wired into the Sphinx toctree.
% Build: pdflatex cobordism_bridge_primer.tex  (run twice for refs/TOC).
\documentclass[11pt]{article}

\usepackage[margin=1in]{geometry}
\usepackage{amsmath, amssymb, amsthm, mathtools, bm, booktabs}
\usepackage[hidelinks]{hyperref}

\theoremstyle{plain}
\newtheorem{theorem}{Theorem}
\theoremstyle{definition}
\newtheorem{definition}{Definition}
\theoremstyle{remark}
\newtheorem{remark}{Remark}

\newcommand{\C}{\mathbb{C}}
\newcommand{\R}{\mathbb{R}}
\newcommand{\Z}{\mathbb{Z}}
\newcommand{\HH}{\mathcal{H}}
\newcommand{\vecop}{\operatorname{vec}}
\newcommand{\Tr}{\operatorname{Tr}}
\newcommand{\rank}{\operatorname{rank}}
\newcommand{\ket}[1]{\lvert #1 \rangle}
\newcommand{\bra}[1]{\langle #1 \rvert}
\newcommand{\braket}[2]{\langle #1 \vert #2 \rangle}
\newcommand{\inner}[2]{\langle #1 , #2 \rangle}
\newcommand{\ip}[3]{\langle #1 \vert #2 \vert #3 \rangle}
\newcommand{\adj}{^{\dagger}}
\newcommand{\geo}{\operatorname{geo}}
\newcommand{\WAB}{W_{\!AB}}

% Self-contained pseudocode box (no algorithm2e / algpseudocode on this box).
\newsavebox{\algoboxsave}
\newenvironment{algobox}[1]
  {\par\medskip\setlength{\fboxsep}{9pt}\noindent
   \begin{lrbox}{\algoboxsave}\begin{minipage}{0.93\textwidth}
   \textbf{Algorithm.}~\textsc{#1}\par\vspace{2pt}\hrule\vspace{4pt}\footnotesize}
  {\end{minipage}\end{lrbox}\fbox{\usebox{\algoboxsave}}\par\medskip}

\title{\textbf{From Spectra to Topology}\\[2pt]
\large The Hodge--Dijkgraaf--Witten bridge: harmonic qubits, holonomy,\\
the Choi bend, and finding geometry by minimizing eigenvector residuals}
\author{tessera --- cobordism realizability programme}
\date{}

\begin{document}
\maketitle

\begin{abstract}
\noindent
This primer assembles the theory behind the realizability bridge. We are after a
single statement of the \emph{state--operation--cobordism correspondence}: a
quantum operation $U$ is realizable as a bulk cobordism $\WAB$ whose value
reproduces the transition amplitude $\ip{\psi_A}{U}{\psi_B}$. There are two
roads to that value. The \emph{spectral} road builds the geometry: states become
harmonic forms of a combinatorial Hodge Laplacian, and we \emph{synthesize} a
bulk whose boundary Laplacian carries the (Choi-bent) operator as an eigenvector.
The \emph{topological} road is the Dijkgraaf--Witten state sum: a metric-free
invariant computed from the cobordism's topology and a group cocycle. The bridge
asks whether the two roads give the same number. We develop, in order: the
combinatorial Hodge Laplacian and its harmonic spaces; the \emph{Hodge qubit};
flat $\Z_2$ holonomy and the Dijkgraaf--Witten partition function; the crucial
distinction between the $b_1$-dimensional harmonic space and the
$2^{\,b_1}$-dimensional holonomy space, and the preparation map between them; the
Choi--Jamio\l kowski ``bend'' that turns an operation into a state; and a
whitepaper-style algorithm that recovers \emph{topology} by driving an
eigenvector residual to zero.
\end{abstract}

\tableofcontents
\bigskip

\section{Overview: two roads to \texorpdfstring{$Z(\WAB)$}{Z(W\_AB)}}

Fix two boundary states $\psi_A, \psi_B$ and an operation
$U:\HH_B\!\to\!\HH_A$. The correspondence claims three things about a bulk
$\WAB$ built ``from'' $U$:
\begin{enumerate}
\item[\textbf{(H1)}] $\WAB = \geo(U)$ is a cobordism (a manifold-with-boundary);
\item[\textbf{(H2)}] $\partial \WAB = \overline{\geo(\psi_B)} \sqcup \geo(\psi_A)$
  --- its boundary is the (outgoing) geometry of $\psi_A$ and the
  orientation-reversed (incoming) geometry of $\psi_B$;
\item[\textbf{(H3)}] $Z(\WAB) = \ip{\psi_A}{U}{\psi_B}$ --- the cobordism's
  invariant equals the transition amplitude.
\end{enumerate}
The subtlety lives in (H3), because ``$Z$'' has two meanings:
\begin{itemize}
\item \textbf{Spectral / geometric $Z$.} Realize $\WAB$ by \emph{synthesizing}
  it: pin its boundary, fill its interior (edge weights and topology) until the
  boundary Hodge--Laplacian eigenvector matches the bent operator. The value is
  read off the geometry; it depends on the metric (edge weights).
\item \textbf{Topological $Z$ (Dijkgraaf--Witten).} Compute the $\Z_2$ state sum
  from the triangulation and a cocycle alone --- no edge weights. The value is a
  topological invariant.
\end{itemize}
The bridge is the claim that these coincide on $\WAB$. The rest of this primer
builds the machinery on both sides and the map between their boundary data.

\section{The combinatorial Hodge Laplacian}

Let $K$ be a (finite, oriented) simplicial complex with weighted edges. Write
$C^k$ for the space of $\C$-valued $k$-cochains (functions on oriented
$k$-simplices, antisymmetric under reorientation). The \emph{coboundary}
$d_k : C^k \to C^{k+1}$ is the signed incidence operator,
\begin{equation}
(d_k \alpha)(\sigma) \;=\; \sum_{\tau \subset \sigma,\ \dim\tau = k}
\,[\sigma : \tau]\,\alpha(\tau),
\end{equation}
with $[\sigma:\tau]\in\{\pm1\}$ the relative orientation. The defining identity
of a chain complex is $d_{k+1} d_k = 0$.

A choice of Hermitian inner products on each $C^k$ (here induced by the edge
weights, see Remark~\ref{rmk:magnitude}) gives adjoints $d_k\adj : C^{k+1}\to
C^k$. The \emph{Hodge Laplacian} in degree $k$ is
\begin{equation}
L_k \;=\; \underbrace{d_{k-1}\,d_{k-1}\adj}_{\text{``down''}}
\;+\; \underbrace{d_k\adj\, d_k}_{\text{``up''}} ,
\qquad L_k = L_k\adj \succeq 0 .
\end{equation}

\begin{remark}[The magnitude convention at $k=0$]\label{rmk:magnitude}
At $k=0$ the down-term vanishes and $L_0 = d_0\adj d_0$ is the graph Laplacian.
With a \emph{Hermitian} weight $A_{ij} = w_{ij}\,e^{\,i\phi_{ij}}$ on the edge
$i\!\to\!j$ (so $A_{ji}=\overline{A_{ij}}$) and the \emph{magnitude} degree
$D_{ii} = \sum_j \lvert w_{ij}\rvert$, the project uses
\begin{equation}
L_0 \;=\; D - A , \qquad D = \operatorname{diag}(D_{ii}),
\end{equation}
a Hermitian operator whose evolution $e^{-iL_0 t}$ is unitary. The phase
$\phi_{ij}$ is a $U(1)$ connection living on the edge; its loop sums are
gauge-invariant fluxes. This is exactly \texttt{cobordism::HodgeLaplacian} (the
class is degree-parameterized: $k=0$ from edge weights, $k\ge 1$ from the
boundary matrices of \texttt{ChainComplex}).
\end{remark}

The central fact is the discrete Hodge theorem.

\begin{definition}[Harmonic cochains]
$\alpha\in C^k$ is \emph{harmonic} iff $L_k\alpha = 0$; equivalently
$d_k\alpha = 0$ and $d_{k-1}\adj\alpha = 0$ (closed and co-closed).
\end{definition}

\begin{theorem}[Discrete Hodge decomposition]
$C^k = \operatorname{im} d_{k-1} \;\oplus\; \ker L_k \;\oplus\;
\operatorname{im} d_k\adj$, and the harmonic space is isomorphic to cohomology,
\begin{equation}
\ker L_k \;\cong\; H^k(K;\,\C), \qquad
\boxed{\;\dim \ker L_k \;=\; b_k\;}
\end{equation}
the $k$-th Betti number. Harmonic representatives are the unique
minimum-norm elements of their cohomology classes.
\end{theorem}

So the \emph{kernel} of $L_k$ is a coordinate-free handle on topology: its
dimension counts $k$-dimensional ``holes,'' and its elements are canonical
(metric-orthogonalized) representatives of those classes. The eigenvectors of
$L_k$ with small nonzero eigenvalue are the ``almost-harmonic'' modes; the
spectrum interpolates between geometry (large eigenvalues, local) and topology
(the kernel, global).

\section{Harmonic forms and the Hodge qubit}\label{sec:qubit}

Take $\Sigma$ a closed oriented surface of genus $g$. Then
$b_1(\Sigma) = 2g$, so $\dim \ker L_1(\Sigma) = 2g$: there are exactly $2g$
independent harmonic $1$-forms, one ``dual'' pair per handle.

\begin{definition}[Hodge qubit]
For the torus $T^2$ ($g=1$) we have $\dim\ker L_1(T^2) = 2$: a two-dimensional
complex vector space. We call it the \emph{Hodge qubit}. A unit harmonic form
\begin{equation}
\psi = \alpha\,\omega_x + \beta\,\omega_y \in \ker L_1(T^2),
\qquad \lvert\alpha\rvert^2 + \lvert\beta\rvert^2 = 1,
\end{equation}
written in the basis of the two harmonic generators $\omega_x,\omega_y$ (the
discrete analogues of $dx, dy$ around the two cycles), \emph{is} the qubit state
$\alpha\ket{0} + \beta\ket{1}$.
\end{definition}

The amplitudes and relative phase of a quantum state are encoded as the
weights and holonomies of a harmonic $1$-form on a surface. Boundary states for
the cobordism are chosen as $\psi_A \in \ker L_1(\Sigma_A)$,
$\psi_B \in \ker L_1(\Sigma_B)$. (The $k=0$ analogue --- a node-amplitude qubit
$\ker L_0$ on a small graph --- is the reduced setting used for the first
realizability oracle; the surface/$k{=}1$ version is the one that matches
Dijkgraaf--Witten, \S\ref{sec:dw}--\S\ref{sec:counting}.)

\section{Holonomy, flat connections, and Dijkgraaf--Witten}\label{sec:dw}

Now switch from continuous harmonic forms to discrete gauge fields. A
\emph{flat $\Z_2$ connection} on a complex assigns $g_e \in \Z_2 = \{0,1\}$ to
each edge so that the holonomy around every triangle is trivial; modulo gauge
transformations these are classified by
\begin{equation}
H^1(\Sigma;\Z_2) \;=\; \operatorname{Hom}\!\big(H_1(\Sigma),\Z_2\big)
\;\cong\; (\Z_2)^{\,b_1(\Sigma)} .
\end{equation}
A class is fixed by its \emph{holonomy} $\pm 1$ around each of the $b_1$
independent cycles, so there are $2^{\,b_1}$ of them.

\begin{definition}[Dijkgraaf--Witten boundary space]
A boundary surface $\Sigma$ carries the Hilbert space
\begin{equation}
Z(\Sigma) \;=\; \C\big[\,H^1(\Sigma;\Z_2)\,\big],
\qquad \dim Z(\Sigma) = 2^{\,b_1(\Sigma)},
\end{equation}
spanned by the holonomy classes (one basis vector per flat connection).
\end{definition}

For a finite group $G$ (here $\Z_2$) and a group $n$-cocycle
$\omega \in H^n(BG; U(1))$, the \emph{Dijkgraaf--Witten state sum} of a closed
$n$-manifold $W$ is
\begin{equation}
Z(W) \;=\; \frac{1}{\lvert G\rvert^{\,\lvert V\rvert}}
\sum_{\substack{g:\,\text{edges}\to G \\ \text{flat}}}
\ \prod_{\text{$n$-simplices } t} \omega(g_t)^{\,\varepsilon_t},
\end{equation}
a number; $\varepsilon_t=\pm1$ is the orientation of the top simplex $t$. For
$n=3$ and $G=\Z_2$, $\omega$ is one of two $3$-cocycles --- the trivial one
($\omega\equiv 1$) and the nontrivial \emph{sign} class
($\omega = (-1)^{abc}$). For a cobordism $W$ with boundary one holds the
boundary connection fixed and sums the same product over interior flat fields;
the result is a vector in $Z(\partial W)$, equivalently (for two boundary
components) a linear map
\begin{equation}
Z(W) : Z(\Sigma_B) \longrightarrow Z(\Sigma_A),
\qquad
\ip{\psi_A}{Z(W)}{\psi_B} \;=\;
\sum_{a,b} \overline{\psi_A[a]}\; Z(W)_{ab}\; \psi_B[b] .
\end{equation}
For the trivial cylinder $\Sigma\times[0,T]$ this map is the identity on
$Z(\Sigma)$. In the codebase this is \texttt{cobordism::DijkgraafWitten}:
\texttt{partitionFunction()} (closed), \texttt{map()} / \texttt{boundaryVector()}
(with boundary), and \texttt{amplitude(}$\psi_A,\psi_B$\texttt{)} for the
sandwiched value above.

\section{Two countings, one preparation map}\label{sec:counting}

We now have \emph{two} boundary spaces attached to the same surface $\Sigma$,
and they are \emph{not} the same:

\begin{center}
\begin{tabular}{@{}lll@{}}
\toprule
 & Hodge qubit & Dijkgraaf--Witten space \\
\midrule
object & harmonic $1$-forms $\ker L_1(\Sigma)$ & holonomy classes $\C[H^1(\Sigma;\Z_2)]$ \\
dimension & $b_1$ \;(continuous, over $\C$) & $2^{\,b_1}$ \;(discrete count) \\
torus $T^2$ & $\C^2$ \;(one qubit) & $\C^4$ \\
role & the spectral state & the state-sum basis \\
\bottomrule
\end{tabular}
\end{center}

\begin{remark}[Keep the countings distinct]
$b_1$ versus $2^{\,b_1}$: continuous harmonic forms give the \emph{qubit
dimension}; flat connections \emph{index the state sum}. They coincide only
trivially. This is stated as a standing convention in the experiment charter
(\S5.2), and the Dijkgraaf--Witten class deliberately keeps the ``$b_1$
qubit'' separate from the ``$2^{\,b_1}$ flat-connection count.''
\end{remark}

The two sides meet through a \emph{preparation map}
\begin{equation}
\mathsf{prep}:\ \ker L_1(\Sigma)\ \longrightarrow\ \C[H^1(\Sigma;\Z_2)],
\qquad b_1 \mapsto 2^{\,b_1},
\end{equation}
which encodes a harmonic (spectral) boundary state as an amplitude over holonomy
classes; \texttt{DijkgraafWitten.amplitude()} already consumes states in the
flat-connection basis that have been prepared this way. Making $\mathsf{prep}$
and its readout first-class and tested is the foundational step of the bridge
work --- it is the only place the spectral and topological descriptions of the
\emph{boundary} must be reconciled. (Everything ``interesting'' then happens in
the \emph{bulk} invariant.)

\section{Choi--Jamio\l kowski: operations as states (the bend)}\label{sec:choi}

To feed an operation into a geometry we must first turn it into a state. The
vehicle is the Choi--Jamio\l kowski isomorphism. Define the \emph{vectorization}
of $U = \sum_{ij} U_{ij}\,\ket{i}\!\bra{j} : \HH_B\!\to\!\HH_A$ by stacking rows,
\begin{equation}
\vecop(U) \;=\; \sum_{ij} U_{ij}\,\ket{i}\otimes\ket{j} \in \HH_A\otimes\HH_B,
\qquad
\vecop\big(\ket{a}\!\bra{b}\big) \;=\; \ket{a}\otimes\overline{\ket{b}} .
\end{equation}
The operator becomes a (generally entangled) bipartite state --- the
\emph{Choi state}. Two facts make it the right bend.

\begin{theorem}[Hilbert--Schmidt / Choi identity]\label{thm:hs}
For the rank-one transition operator $U_T = \ket{\psi_A}\!\bra{\psi_B}$,
\begin{equation}
\ip{\psi_A}{U}{\psi_B} \;=\; \Tr\!\big(U_T\adj\,U\big)
\;=\; \inner{\vecop(U_T)}{\vecop(U)} .
\end{equation}
\end{theorem}

\noindent
That is: the physical amplitude $\ip{\psi_A}{U}{\psi_B}$ equals an
\emph{inner product of Choi states}. So if a cobordism faithfully carries
$\vecop(U)$ on its output boundary, contracting that boundary state against the
test functional $\vecop(U_T)$ returns the amplitude --- which is exactly how the
spectral side reads off ``$Z$.'' (This is \texttt{quantum::ChoiJamiolkowski}:
\texttt{vectorize}, \texttt{transitionOperator}, \texttt{transitionAmplitude}.)

\begin{theorem}[Rank $=$ Schmidt rank $=$ connectivity]
The Schmidt rank of $\vecop(U)$ equals $\rank(U)$ (the number of nonzero singular
values of $U$). A separable Choi state ($\rank 1$) corresponds to a rank-one
$U_T=\ket{\psi_A}\!\bra{\psi_B}$ and a \emph{disconnected} cup/cap geometry;
higher Schmidt rank is genuine entanglement and a \emph{connected} bulk.
\end{theorem}

\section{Finding topology by minimizing eigenvector residuals}\label{sec:residual}

We can now state the inverse problem precisely. Given a target boundary state
$\psi$ (the bent operator, $\psi=\vecop(U)/\lVert\vecop(U)\rVert$), \emph{find a
complex $W$ whose Hodge Laplacian has $\psi$ as an eigenvector}. The
eigenvalue is a free output, so we use the eigenvalue-agnostic
\emph{residual}
\begin{equation}
\boxed{\;
r(W) \;=\; \big\lVert\, \big(I - \psi\psi\adj\big)\,L(W)\,\psi \,\big\rVert^2
\;=\; \big\lVert\, L\psi - (\psi\adj L \psi)\,\psi \,\big\rVert^2 .
\;}
\end{equation}
$r(W)$ is the squared norm of the component of $L\psi$ orthogonal to $\psi$; it
vanishes \emph{iff} $L\psi = \lambda\psi$ with $\lambda = \psi\adj L\psi$ the
Rayleigh quotient. Minimizing $r$ over the degrees of freedom of $W$ is thus a
search for a geometry that \emph{realizes} $\psi$ spectrally.

\paragraph{Fixed-boundary fill.} For realizability the boundary is
\emph{pinned}: $\partial W$ is the synthesized $\geo(\psi_A),\geo(\psi_B)$ and the
output surface, held fixed. The free parameters $\theta$ are the \emph{interior}
edge weights/phases and --- through boundary-fixed Pachner ``add'' moves (coning)
--- the interior \emph{topology}. This is the decisive twist: with the ends fixed
the parameters do \emph{not} outrun the constraints, so realizability becomes a
genuine test rather than a foregone conclusion.

\paragraph{Verdict and certificate.} A target is \textbf{realizable} iff $r$ can
be driven below $\epsilon$; it is \textbf{obstructed} iff $r$ \emph{floors} at a
positive value --- a residual lower bound under the fixed-boundary constraint.
The floor is the non-existence \emph{certificate} (the analogue of the
two-vertex floor $w_{\min}^2(\lvert c_0\rvert^2-\lvert c_1\rvert^2)^2$), obtained
without exhausting triangulations.

\begin{algobox}{DecideRealizability$(U,\,d_A,\,d_B;\ \epsilon,\,R,\,c_{\max})$}
\textit{Input:} operation $U:\HH_B\!\to\!\HH_A$; tolerance $\epsilon$; restart
count $R$; cone budget $c_{\max}$.\\
\textit{Output:} \textsc{realizable} with witness $W$, or \textsc{obstructed}
with floor $f$.\\[3pt]
\rule{\linewidth}{0.4pt}\\[4pt]
$\psi \leftarrow \vecop(U)\,/\,\lVert\vecop(U)\rVert$
   \hfill\textit{// bend: Choi--Jamio\l kowski (\S\ref{sec:choi})}\\[1pt]
$W \leftarrow$ seed bulk;\ \ \textsc{Pin}$(\partial W)$
   \hfill\textit{// fixed ends; only interior edges are writable}\\[1pt]
$f \leftarrow +\infty$\\[1pt]
\textbf{for } $c = 0,1,\dots,c_{\max}$ \textbf{ do}\\[1pt]
\hspace*{1.6em} define $r_W(\theta) = \lVert (I-\psi\psi\adj)\,L(W;\theta)\,\psi\rVert^2$
   \hfill\textit{// residual over interior params $\theta$}\\[1pt]
\hspace*{1.6em} $(\theta^\star, r^\star) \leftarrow
   \textsc{MultiRestartLM}\big(r_W,\ R\big)$
   \hfill\textit{// Levenberg--Marquardt $\times R$ seeds}\\[1pt]
\hspace*{1.6em} \textbf{if } $r^\star < \epsilon$ \textbf{ then return }
   \textsc{realizable},\ $W(\theta^\star)$
   \hfill\textit{// witness carries $\psi$ as an eigenvector}\\[1pt]
\hspace*{1.6em} $f \leftarrow \min(f,\ r^\star)$\\[1pt]
\hspace*{1.6em} $W \leftarrow \textsc{ConeInterior}(W)$
   \hfill\textit{// boundary-fixed Pachner add: +1 interior vertex}\\[1pt]
\textbf{return } \textsc{obstructed},\ \text{floor } f
   \hfill\textit{// $f$ bounded away from $0$ certifies non-existence}
\end{algobox}

\noindent
In the codebase: \texttt{ChoiJamiolkowski} does the bend;
\texttt{EigenstateSynthesis} evaluates $r_W$ and the fixed-boundary fill;
\texttt{LevenbergMarquardt} is the multi-restart solver;
\texttt{RealizabilityOracle.decide()} is the loop above; the witness $W$ is
returned as a \texttt{Spacetime}, and \texttt{getBoundary()} recovers
$\partial W$.

\section{The bridge}\label{sec:bridge}

Putting the pieces together on a single synthesized $\WAB$:
\begin{itemize}
\item the spectral road (\S\ref{sec:residual}) produces $\WAB$ whose
  boundary harmonic carries $\vecop(U)$, so by Theorem~\ref{thm:hs} the
  \emph{spectral} $Z_{\mathrm{spec}}(\WAB) = \ip{\psi_A}{U}{\psi_B}$;
\item the topological road (\S\ref{sec:dw}) computes
  $Z_{\mathrm{DW}}(\WAB) = \ip{\psi_A}{Z(\WAB)}{\psi_B}$ from the triangulation
  and cocycle alone, with boundary states obtained through the preparation map
  of \S\ref{sec:counting}.
\end{itemize}
The bridge is the assertion
\begin{equation}
Z_{\mathrm{DW}}(\WAB) \;\overset{?}{=}\; \ip{\psi_A}{U}{\psi_B}
\;\overset{?}{=}\; Z_{\mathrm{spec}}(\WAB).
\end{equation}

\begin{remark}[A discrete shadow]
The $\Z_2$ Dijkgraaf--Witten functor sends cobordisms to a \emph{discrete}
family of maps (the Frobenius algebra $\C[\Z_2]$ generated by pair-of-pants,
cap/cup, cylinder, twisted by the cocycle). A generic $U$ is \emph{not}
$Z_{\mathrm{DW}}$ of any $\Z_2$ cobordism. The honest expectation, then, is that
the bridge holds on the discrete \emph{DW-representable} subset, and that the
spectral oracle --- which realizes a \emph{continuum} of operators --- strictly
extends the topological theory. In that picture the Dijkgraaf--Witten invariant
is the \emph{quantized shadow} of the spectral $Z$: equal where both are defined,
with the geometry seeing strictly more than the finite-group state sum.
\end{remark}

\paragraph{Where this lives.} \texttt{HodgeLaplacian} (\S2--\S3),
\texttt{DijkgraafWitten} (\S4--\S5), \texttt{ChoiJamiolkowski} (\S6), and
\texttt{RealizabilityOracle}/\texttt{EigenstateSynthesis}/\texttt{LevenbergMarquardt}
(\S7) are in place; the preparation map of \S\ref{sec:counting}, the $k{=}1$
3-manifold synthesis, and the agreement experiment of \S\ref{sec:bridge} are the
three steps of the bridge milestone.

\end{document}
